\scanpage{81}
% archive foliation: 80

est un \uline{homomorphisme d'anneaux} (c'est en effet le Frobenius habituel\deleted{\unsure{)}},
\deleted{et \unsure{on a} $V,D$ sont des \uline{morphismes additifs} (}$V$ est la Verschiebung habituelle) et on a
\[
FV=p.\,\mathrm{id},\quad VF=p.\mathrm{id}\qquad(\text{\uline{en car.} }p).
\]
On introduit aussi le schéma en groupes en
\[
\underrightarrow{W}=\varinjlim_n\underline{W}_n,\quad\text{morphismes de transition du type }D.
\]
Ainsi du point de vue schémas,on a
\[
\underrightarrow{W}=\underline{E}^{(\underline{N})}.
\]
\uline{En car.} $p$, on a commutat\deleted{t}ivité dans
\diagram{\begin{tikzcd}
\underline{W}_n \arrow[r,"T"] & \underline{W}_{n+1}\\
\underline{W}_n \arrow[u,"F"] \arrow[r,"p"] & \underline{W}_{n+1} \arrow[u,"\mathrm{id}"']
\end{tikzcd}}
\deleted{d'où}\inserted{et} commutativité dans
\diagram{\begin{tikzcd}
\underline{W}_n \arrow[r,"T"] \arrow[d,"F^{n-1}"'] & \underline{W}_{n+1} \arrow[d,"F^{n\deleted{\unsure{-1|+1}}}"]\\
\underline{W}_n \arrow[r,"p"] & \underline{W}_{n+1}
\end{tikzcd}},
d'où un homomorphisme canonique
\[
u:\underrightarrow{W}_{\inserted{\underline{F}_p}}\longrightarrow
\varinjlim(\underline{W}_{n\inserted{\underline{F}_p}},p),
\]
qui \deleted{est un} induit un isomorphisme sur les points à valeurs dans un \deleted{corps}\inserted{anneau} parfait.

Soit maintenant $k$ un corps parfait de car. $p$. Alors
\[
W=\underline{W}(k)=\varprojlim W_n\qquad(W_n=\underline{W}_n(k))
\]
est un anneau de valuation discrète de corps résiduel \deleted{\unsure{k, dit}} $W_1=k$, d\deleted{e}'idéal maximal engendré par $p$, et de corps des fractions
\[
K=W[1/p]
\]
de caractéristique nulle. L'homomorphisme $u$ précédent induit un isomorphisme de groupes
\[
u:\underrightarrow{W}(k)\xrightarrow{\sim}K/W.
\]
On va introduire sur $\unsure{\underline{W}|\underrightarrow{W}}_k$ une structure de schéma en groupes avec
